\documentclass{ximera}
\title{Simple Harmonic Motion, Summary}

\newcommand{\pskip}{\vskip 0.1 in}

\begin{document}
\begin{abstract}
Simple harmonic motion, a summary of the key ideas.
\end{abstract}
\maketitle

\section{Mass on a Spring}

We'll start the class with one of its most important topics, simple harmonic motion. The classic example starts with a mass attached to the end of a spring. Stretch the spring beyond its relaxed length and relase the mass. Then then mass, assumed free to slide without friction on a horizontal surface, oscillates between its extremes in emphasis \emph{simple harmonic motion}.

Such oscillatory motion is common and one primary aim of this class is to describe the motion algebraically. While we do not yet have the tools to do this, we can still describe the motion geometrically. And it's this geometric description as illustrated in the following example that forms the basis not only for simple harmonic motion but for much of this course.

{\bf Note:} If you already had trigonometry, forget what you know for now. In particular, forget about the sine and cosine functions. The purpose of these examples is to think geometrically, not algebraically.



\begin{example}  \label{Ex:LLL}
A block $M$ at the end of a spring oscillates in simple harmonic motion on a frictionless surface as shown below. The block is stretched two meters from beyond its equilibrium position and then released from rest (at $s=2$ m), completing one oscillation about equilibrium ($s=0$) every $5$ seconds. 

\begin{onlineOnly}
    \begin{center}
\desmos{dcba538898}{900}{600}
\end{center}
\end{onlineOnly}

\href{https://www.desmos.com/calculator/dcba538898}{163: Simple Harmonic Motion 23}

\begin{enumerate}
\item Sketch by hand a graph of the function
\[
      s = f(t) \, , \, t\geq 0,
\]
that expresses the displacement of the block (measured in meters) from equilibirium in terms of the number of seconds since the block was released. The displacement is positive when the block is to the right of equilibrium, negative when it is to the left. Be sure to label the coordinate axes with the appropriate variable names and units.

\emph{Note:} The slider $u$ in Line 2 is another name for the time $t$ (measured in seconds) since the mass was released. You can stop the motion and then drag the slider to help approximate coordinates of points on the graph of the function $f$.

\item Activate the folder \emph{Position function graph} in Line 7 by clicking the open circle on Line 7 to check your graph.

\item To say the block oscillates in \emph{simple harmonic motion} means something much more specific than that it simply oscillates back and forth. To understand exactly what simple harmonic motion means, activate the folder \emph{circle} on Line 3 and resume the motion.

The key idea things to observe are these:

\begin{enumerate}

\item The point $P$ moves around the circle at a \emph{at a constant speed}.

\item This circle is centered at the equilibrium position $=0$ and has radius equal to the maximimum displacement of the block from equilibrium (two meters in this example).

\item The block $M$ stays directly above or below the point $P$ at all times.

\end{enumerate}

%\emph{Uniform circular motion drives simple harmonic motion along a diameter of the circle.}



%\item Use the animation above to apprimate the \emph{displacement} of the mass from the origin (ie. its $x$-coordinate) at times $t=0, 1, 2, 4$ seconds. The displacement (ie. $x$-coordinate) is positive when the mass is to the right of the origin, negative to the left. 

%\item Use the animation to approximate the displacement from the origin at time $t=26.5$ seconds.

%\item Use the animation to approximate the first four times when the mass is $1$ meter to the left of the origin.

%\item Use the animation to approximate when the mass is $1$ meter to the left of the origin for the $103$rd time.
\end{enumerate}
\end{example}

The key takeaway from the first example is that \emph{projecting a constant speed motion around a circle onto a diameter results in simple harmonic motion along that diameter}.  

The next example takes the first step toward an algebraic description of simple harmonic motion by replacing the animation with a protractor.



\begin{example} \label{ExL9df93r3}
The scenario here is exactly the same as in Example 1. But the animation is different. Now time is surpressed and replaced with the polar angle $v$, measured counterclockwise (in \emph{revolutions}) from the positve $s$-axis.

\begin{onlineOnly}
    \begin{center}
\desmos{um74xweymf}{900}{600}
\end{center}
\end{onlineOnly}

\href{https://www.desmos.com/calculator/um74xweymf}{142: Simple Harmonic Motion 24}

\begin{enumerate}
\item Find a function 
\[
     \theta = a(t) \, , \, t\geq 0,
\]
that expresses the polar angle (in revolutions, measured counterclockwise from the positive $s$-axis) of the driving point $P$ in terms of the number of seconds since the mass was released. Include units in your computation.

The driving point $P$ rotates about the center of its circular path at the constant rate of 
\[
     \omega = \frac{\Delta \theta}{\Delta t}  = \answer{\frac{1}{5}} \frac{\text{rev}}{\text{sec}}.
\]

Since $\theta=0$ at time $t=0$, the polar angle function is
\[
   \theta = a(t) =  \answer{t/5} \, , \, t \geq 0.
\]

\item Use your polar angle function from part (a) and the protractor above (where the ticks are space at intervals of $0.01$ revolutions) to approximate answers to the following questions.

\begin{enumerate}
\item Find the position of the block $4$ seconds after being released.

\item Find the position of the block $24.5$ seconds after being released.

\item Find the first four times when the block is $0.5$ meters to the left of equilibrium.

\item When is the block is $0.5$ meters to the left of equilibrium for the $103$rd time?
\end{enumerate}
\end{enumerate}

\begin{explanation}
\begin{enumerate}
\item First we use the polar angle function to find the polar angle of the driving point at time $t=4$. This angle is
\[
 \theta = a(4) = \left( \frac{1}{5}  \frac{\text{rev}}{\text{sec}}  \right) (4 \text{ sec}) = \frac{4}{5} \text{rev}.
\]
Now use the protractor and drag the slider in $v$ (another name for $\theta$) in Line 2 to $v=0.8$, so that the polar angle of $P$ is $\theta = 8/10$ rev. Then approximate the $s$-coordinate of the block as $s\sim 0.6$ cm. So the block is approximately $0.6$ meters to the right of its equilibrium position four seconds after being released.
\end{enumerate} 
\end{explanation}

\end{example}



We do not necessarily want to rely on a protractor tied specifically to the question at hand. The next example requires scaling between a universal protractor to a specific simple harmonic motion. 



\begin{example}  \label{Ex:FEfe344trgfd45}

Suppose for this problem that the earth is a ball with uniform density of radius 4000 miles. Now imagine drilling a straight tunnel through the earth from the North Pole to the South Pole. A rock dropped from rest at the north pole falling through the tunnel would then oscillate in simple harmonic motion between the poles and return to the north pole every 84  minutes as illustrated below. %The rock reaches the north pole at 12:23pm.

\begin{onlineOnly}
    \begin{center}
\desmos{z5ya8iy2e4}{900}{600}
\end{center}
\end{onlineOnly}

\href{https://www.desmos.com/calculator/z5ya8iy2e4}{142: Tunnel Through the Earth }

Suppose the oscillating rock above passes the north pole at 12:23pm.

\begin{enumerate}

\item Find a function
\[
   \theta = a(t) \, , \, 0\leq t \leq 24,
\]
that expresses the polar angle of the point $Q$ (measured counterclockwise in revolutions from the north pole) in terms of the number of minutes past noon. Use point-slope.

The polar angle function is
\[
   \theta = a(t) = \answer{\frac{1}{84}}\left( t - \answer{23} \right).
\]

\item Find a function
\[
   h = f(x) \, , \, 0\leq x \leq 8,
\]
that expresses the distance of the rock (in thousands of miles) from the south pole in terms of the $x$-coordinate of the driving point $P$ on a protractor of radius $5$ cm. Use point-slope.

The function is
\[
     h = f(x) = \answer{\frac{4}{5}} \left(\answer{x+5}\right) \, , \, -5 \leq x \leq 5.
\]

\item Use your functions from parts (a) and (b), the protractor of 5 cm in this example, and four-function arithmetic to answer the following questions. Activate the \emph{Protractor} folder in Line 15 of the worksheet above if you need help. 
\begin{enumerate}

\item Estimate the distance of the rock from the South Pole at 12:40pm.

\item Estimate the first two times after 12:00pm when the rock is $1000$ miles from the South Pole.

\item During one oscillation, estimate the fraction of the time is the rock at least $1200$ miles from the South Pole.
\end{enumerate}

\end{enumerate}
\end{example}

\begin{example} \label{Ex:Lfe3r340943210}
On the first day of fall, every location on earth gets 12 hours of daylight/day. On the first day of winter, Shoreline gets $8$ hours of daylight/day.

To make things simple, suppose each month has 30 days and that the winter solstice occurs 90 days after the fall equinox.

\begin{enumerate}
\item The time from sunrise to sunset on the second day of fall is a bit less than 12 hours. Assume the number of hours of daylight/day in Shoreline is a piecewise linear function and approximate the difference in the number of \emph{minutes} of daylight/day in Shoreline between these two days..

\item Assume instead that over the course of the year the number of hours of daylight/day in Seattle oscillates in simple harmonic motion about its mean of 12 hours/day and use the protractor in Example 3 to approximate the same difference.

If you have trouble, you can try using the protractor below instead.

\begin{onlineOnly}
    \begin{center}
\desmos{dhzo6hhjw1}{900}{600}
\end{center}
\end{onlineOnly}

\href{https://www.desmos.com/calculator/dhzo6hhjw1}{142: Circular Interpolation }



\item What fraction of the fall need pass by for the number of daylight hours/day in Shoreline to decrease halfway (ie. to 10 hours of daylight/day) from its mean of 12 hours of daylight/day on the fall equinox to its minimum of 8 hours of daylight/day on the winter solstice? How would this fraction change for Fairbanks, AK, where they get only $4$ hours of daylight/day on the winter solstice? Use the protractor of Example 2 or the protractor above to approximate these fractions.

\end{enumerate}
\end{example}

\begin{example}
You ride a ferris wheel that rotates at a constant rate. What fraction of the time do you spend at a height at least halfway between the wheel's center and its high point? Use the protractor of Example 3 to approximate this fraction.
\end{example}


\end{document}
